\section*{Abstract}
\addcontentsline{toc}{section}{Abstract (EN)}

\textbf{Title:} Development of a Telegram Bot Information System for
Public Transport in Fergana City.

\medskip

Public transport in Fergana, a city of about 300{,}000 inhabitants in
Uzbekistan, has no official digital information system. Residents and
visitors depend on word of mouth, paper signs, or general map services
such as Yandex Maps, which do not provide full data on local bus routes.
This creates daily problems: people miss buses, take wrong routes, or
spend extra time at stops.

The aim of the thesis is to design and develop a Telegram bot
information system that provides public transport information for
Fergana city in a simple and accessible way. Telegram was chosen because
it is one of the most popular messengers in Uzbekistan and works well
on low-end smartphones with limited internet.

The bot was developed in Python using the \texttt{aiogram} framework
and SQLite as the database. Route and stop data were collected from
Yandex Maps and stored in a local database. The bot supports six bus
routes and the following main functions: showing route stops, searching
a route between two stops, providing taxi information, and collecting
user suggestions. The system was deployed on the PythonAnywhere cloud
platform for continuous operation.

The bot was tested with a small group of users in Fergana. Results show
that the system is easy to use, responds quickly, and covers the most
common transport questions. The thesis also discusses future
improvements, including integration with the Grok AI model for natural
language queries.

The work confirms that a Telegram bot is a practical and low-cost
solution for cities where official transport applications do not exist.

\medskip

\noindent\textbf{Volume:} XX pages, YY tables, ZZ figures,
zx references, xz appendices.
\newpage
